% ======================================================================
% SECTION 5: LIMITATIONS AND FUTURE WORK
% Five thematic subsections expanded from the v0.8.0 bracketed draft at
% 2030-pdac-1min/paper/draft-paper/sections/limitations_future.tex.
% Location: 2030-pdac-1min/paper/full-paper/sections/limitations_future.tex
% ======================================================================

\subsection{Approximations vs Generated vs Executed Accounting}
\label{sec:limits-accounting}

The five phase workflow makes the approximation budget explicit
at every stage. The v0.5.0 instruction set at
\texttt{2030-pdac-1min/paper/instructions/README.md} pre commits a
per iteration budget of 980 KB across pyramid levels L1 (50 ms
aggregates), L2 (1 s aggregates), L3 (phase aggregates), L4
(anastomosis events), and the event log; the 33.4 MB committed
total across 32 iterations plus 2.0 MB fixed overhead is a
deliberate L0 to L4 down sampling.\ The L0 raw at 412 MB per
iteration and 13.2 GB across 32 iterations is archived to Zenodo
per \texttt{zenodo_archive_protocol.md} and is never committed to
Git.\ The per pyramid level budget is detailed in the file
\texttt{2030-pdac-1min/paper/instructions/file_size_pyramid_1min.md}
and the catalog of GBM 7 approximations addressed by the PDAC
variant is detailed in
\texttt{2030-pdac-1min/paper/instructions/gbm_errors_addressed.md}.

The v0.6.0 codegen tree at
\texttt{2030-pdac-1min/paper/codegen/README.md} and the v0.7.0
execution tree at
\texttt{2030-pdac-1min/paper/execution/README.md} plus
\texttt{PROCESS.md} document what was authored versus what was
approximated and what was actually run versus what was deferred.
The 15 entry cross commit reference matrices at
\texttt{2030-pdac-1min/paper/codegen/CROSS_REFERENCES.md} and
\texttt{execution/CROSS_REFERENCES.md} tie each codegen module to
the executed run output. The smoke test status at
\texttt{execution/tests/test_status.txt} records 10 of 13 smoke
tests passing with 3 known v0.6.0 codegen discrepancies documented
for follow up; closing those discrepancies is the first item in
the data approximation reduction roadmap (see \S
\ref{sec:limits-data}). The Zenodo L0 raw deposition pointer
family at
\texttt{2030-pdac-1min/paper/execution/zenodo/\{deposition_summary.txt, manifest.json, run_00000_L0_raw.zenodo_pointer.json, README.md\}}
holds a deposition size estimate of 13.21 GB pending live upload
once a Zenodo API token is provisioned.

\begin{table}[H]
\caption{4 phase accounting: approximated vs generated vs executed across the v0.5.0 to v0.9.0 workflow.}
\label{tab:limits-accounting}
\centering
\begin{tabular}{>{\raggedright\arraybackslash}p{2.7cm} >{\raggedright\arraybackslash}p{3.2cm} >{\raggedright\arraybackslash}p{2.9cm} >{\raggedright\arraybackslash}p{6cm}}
\toprule
\textbf{Phase} & \textbf{Artifacts} & \textbf{Size} & \textbf{Top approximation flag} \\
\midrule
  Instructions (v0.5.0) & 26 files & 175 KB & 980 KB per iteration commit budget  \\
  Code gen (v0.6.0) & 15 src modules + cfg + schemas & 1.5 MB & Per arm DH parameters are reference values \\
  Execution (v0.7.0) & Sixteen subdirectories & 33 MB & L0 raw 13.2 GB stays on Zenodo \\
  Draft template (v0.8.0) & 8 .tex + bib + README & 30 KB + 1 zip & Bracketed instructions, not yet processed \\
  Full paper (v0.9.0) & 8 .tex + bib + README + zip & 60 KB + 1 zip & Brackets resolved into prose; final pass complete \\
\bottomrule
\end{tabular}
\end{table}

\subsection{60 Minute vs 1 Minute Structural Time Delta}
\label{sec:limits-time}

The 60 minute versus 1 minute structural time delta pattern
inherited from the v0.4.0 GBM paper at
\texttt{2030-gbm-1min/paper/full-paper/final-paper/sections/limitations_future.tex}
extends to the 4 to 8 hour versus 1 minute Whipple compression in
this PDAC variant surgery.\ The Whipple structural time baseline from
\texttt{2030-pdac-1min/paper/instructions/README.md} is 4 to 8
hours of in room time; the 60 second target is a 240x to 480x
compression that no current commercial robot achieves.\ This delta
is a limitation rather than a deliverable: tissue creep,
anastomosis healing, pharmacokinetics, and clinical decision
intervals do not compress with the surgical motion compression.\
The per iteration Daraxonrasib cancer drug serum concentration trajectory at
\texttt{2030-pdac-1min/paper/execution/daraxonrasib/perioperative_trajectory.csv}
is computed at the simulation time scale, not the biological time
scale.

The honest accounting is that the 60 second simulation is a
computational checkpoint, not a clinical readiness claim. The
parallel GBM checkpoint at \cite{kawchak_2026_20113157} is the
direct cross study anchor, and ICH E6(R3) Good Clinical Practice
\cite{ich-e6r3} sets the eventual clinical readiness floor.
Table \ref{tab:limits-time} lists the 5 axes of the 60 minute
versus 1 minute delta extended to the Whipple compression.

\clearpage

\begin{table}[H]
\caption{60 minute vs 1 minute structural delta extended to the 4 to 8 hour vs 1 minute Whipple.}
\label{tab:limits-time}
\centering
\begin{tabular}{>{\raggedright\arraybackslash}p{3.3cm} >{\raggedright\arraybackslash}p{3.4cm} >{\raggedright\arraybackslash}p{3.7cm} >{\raggedright\arraybackslash}p{4.4cm}}
\toprule
\textbf{Axis} & \textbf{Human baseline} & \textbf{1 min target} & \textbf{Compression rationale} \\
\midrule
  Surgical motion & 4 to 8 hours & 60 seconds & 240x to 480x; New 1.0 1,200 mm/s tip \\
  Anastomosis healing & weeks & not modeled in 60 s & Out of scope; downstream PK model needed \\
  Tissue creep & hours to days & not modeled in 60 s & Out of scope; downstream FE model needed \\
  Daraxonrasib serum kinetics & hours to days & sampled at 60 s & PK clock decoupled from surgical clock \\
  Clinical decision intervals & minutes to hours & not modeled in 60 s & Out of scope; downstream IRB protocol needed \\
\bottomrule
\end{tabular}
\end{table}

\subsection{Data Approximation Reduction Roadmap}
\label{sec:limits-data}

Four data approximations live across the five phases and each
closes with a discrete next step. First, the 980 KB per iteration
committed budget at
\texttt{2030-pdac-1min/paper/instructions/file_size_pyramid_1min.md}
is a deliberate L1 to L4 down sampling; the full L0 raw at 412 MB
per iteration is approximated by the Zenodo deposition pointer
under \texttt{2030-pdac-1min/paper/execution/zenodo/}. Provisioning
a Zenodo API token, performing the live upload, and patching the
pointer JSON files in place closes the gap.\ Second, the synthetic
patient PAT-PDAC-0001 information located at
\texttt{2030-pdac-1min/paper/instructions/pdac_context_1min.md} is
one synthetic patient; the downstream cohort scale up under
TRIPOD+AI \cite{Collins2024TRIPODAI} to 100 synthetic PDAC patients
with covariate variation (age, ECOG, CA 19-9, vessel angle, gland
texture) reduces this gap.

Third, the DH parameters at
\texttt{2030-pdac-1min/paper/codegen/config/kinematics_8arm.yaml}
are reference values for the hypothetical PancreSpeed 1.0; bench
measurements on a physical prototype arm reduce this gap, and a
bench prototype of one arm is item 2 in the future work table
(\S \ref{sec:limits-future}). Fourth, the smoke test 10 of 13
pass at \texttt{2030-pdac-1min/paper/execution/tests/test_status.txt}
has 3 known v0.6.0 codegen discrepancies (a vascular gate
boundary off by one tick on the no_fly to soft_warning transition,
a Daraxonrasib advisory enum case sensitivity mismatch, and a
metrics weight serialization order mismatch); closing the 3
discrepancies reduces this gap. The 4 next steps are ordered by
expected impact on the composite score variance (currently std
1.225 around mean 93.298): cohort scale up first because it adds
real variance, smoke test discrepancy closure second because it
removes spurious variance, Zenodo upload third because it does
not affect variance but unlocks reproducibility, and DH parameter
bench measurement fourth because it depends on a physical arm.

\subsection{Surgery Time and Time Resolution Roadmap}
\label{sec:limits-resolution}

The current 60 second surgery time and 100 microsecond force
resolution settings define a tractable but coarse working point. Two
reduction tracks open on the surgery time axis.\ Track A reduces
the simulation surgery time from 60 seconds to 30 seconds; the
PancreSpeed 1.0 spec at
\texttt{2030-pdac-1min/paper/instructions/robot_specification_pancrespeed.md}
would need to support 2,400 mm/s peak tip velocity (2x the current
1,200 mm/s) and the per phase budgets in
\texttt{pdac_context_1min.md} would need to halve.\ Track B
increases the simulation surgery time from 60 seconds to 120
seconds; the biological realism of tissue creep, anastomosis
healing, and pharmacokinetics improves, and the advisory layer
can carry more inter phase deliberation cycles.

Two refinement tracks open on the time resolution axis. Track C
refines to 10 microseconds (1 MHz force sampling and 100 kHz
command); this requires 10x more storage at the L0 raw layer
(132 GB per iteration), which forces a hierarchical storage tier
above Zenodo. Track D coarsens to 1 millisecond (1 kHz force and
1 kHz command); this removes the exceptional processing feat at
\texttt{2030-pdac-1min/paper/execution/sensors/sensor_sample_8arm.jsonl}
but enables a 100x larger iteration sweep (3,200 iterations under
the same per iteration budget). The recommendation is to pursue
Track B (longer surgery, more deliberation) and Track D (coarser
resolution, larger sweep) in parallel because they are compatible
and they each tighten the composite score CI without requiring
new hardware; Track A and Track C are deferred to a later PR
because they require new hardware envelopes.

\subsection{10 Concrete Future Work Deliverables}
\label{sec:limits-future}

Ten concrete future work deliverables follow from the data
approximation reduction roadmap in \S \ref{sec:limits-data} and
the surgery time roadmap in \S \ref{sec:limits-resolution}. The
list is patterned after the v0.4.0 GBM future work table at
\texttt{2030-gbm-1min/paper/full-paper/final-paper/sections/limitations_future.tex}
\cite{kawchak_2026_20113157} and is extended with the PDAC
specific items that the present work surfaces. Each deliverable
is sized to fit a single follow on pull request and to land
inside a 6 to 12 month window. Table \ref{tab:limits-future}
lists the 10 deliverables and the gap that each closes.

\begin{table}[H]
\caption{10 concrete future work deliverables for the 60 second 8 arm Whipple plus Daraxonrasib pairing.}
\label{tab:limits-future}
\centering
\begin{tabular}{>{\raggedright\arraybackslash}p{0.7cm} >{\raggedright\arraybackslash}p{8cm} >{\raggedright\arraybackslash}p{6.7cm}}
\toprule
\textbf{\#} & \textbf{Deliverable} & \textbf{Gap closed} \\
\midrule
  1 & Cohort scale up to 100 synthetic PDAC patients & Single patient PAT-PDAC-0001 limit \\
  2 & Bench prototype of one PancreSpeed 1.0 arm & 5x to 500x SOTA hardware gap \\
  3 & Porcine cadaver Whipple at 5 min target & Tissue creep + anastomosis healing \\
  4 & Track B: 120 second surgery time variant & Inter phase deliberation cycles \\
  5 & Track D: 1 ms time resolution + 3,200 iter sweep & Larger composite score CI tightening \\
  6 & Daraxonrasib PK model with patient covariates & Decoupled pharmacology clock realism \\
  7 & Multi vendor handoff PR, cross licensed data & Simulation vs simulation gap \\
  8 & FDA Q-Sub for the on premises LLM advisory & SaMD pre submission gap \\
  9 & TRIPOD+AI plus CREMLS reporting checklist & Reporting standard compliance gap \\
  10 & Federation across 5 academic centers & Single site evaluation gap \\
\bottomrule
\end{tabular}
\end{table}
